\documentclass{beamer}
\usepackage[utf8]{inputenc}
\usepackage[T1]{fontenc}
\usepackage[spanish, es-tabla]{babel}
\usepackage{helvet}
\usepackage{amsmath, amssymb, amsthm}
\usepackage{graphicx} 
\usetheme{UNAM} 

\title{Procesos Estocásticos}
\subtitle{Generador Infinitesimal en Tiempo Continuo}
\author{Ricardo López Ramírez}
\date{\today}
\institute{\url{423092556@pcpuma.acatlan.unam.com.mx}}

\begin{document}
\begin{frame}[plain,t]
\titlepage
\end{frame}

\begin{frame}
\frametitle{Índice}
\tableofcontents
\end{frame}


\section{Construcción Matricial}

\begin{frame}
\frametitle{Construcción Matricial}
\framesubtitle{Definición de las tasas de transición}

Sea $(X_t)_{t \ge 0}$ una Cadena de Markov en Tiempo Continuo. El sistema se rige por \textbf{tasas instantáneas de transición}, definidas por el límite de $p_{ij}(h) = P(X_{t+h} = j \mid X_t = i)$ cuando $h \to 0^+$:

\begin{block}{Definición de Tasas}
\small  
\begin{itemize}
    \item \textbf{Tasa de salto hacia $j \neq i$:}
    \[ q_{ij} = \lim_{h \to 0^+} \frac{p_{ij}(h)}{h} \implies p_{ij}(h) = q_{ij}h + o(h) \]
    
    \item \textbf{Tasa de salida del estado $i$:}
    \[ q_i = \lim_{h \to 0^+} \frac{1 - p_{ii}(h)}{h} \implies p_{ii}(h) = 1 - q_ih + o(h) \]
\end{itemize}
\end{block}

\end{frame}


\begin{frame}
\frametitle{Construcción Matricial}
\framesubtitle{Propiedades de la Matriz Generadora $Q$}

Agrupando estas tasas formamos la \textbf{matriz generadora infinitesimal} $Q$, con entradas extradiagonales $q_{ij}$ y diagonal $q_{ii} = -q_i$.

Para que $Q$ sea un generador válido, debe cumplir:

\end{frame}

\begin{frame}
\frametitle{Construcción Matricial}
\framesubtitle{Interpretación Probabilística}

La matriz $Q$ nos permite simular la trayectoria del proceso separando el \textit{cuándo} y el \textit{a dónde} salta el sistema.

\begin{theorem}[Construcción mediante saltos]
Si $Q$ es conservativa y estable, la dinámica en el estado $i$ equivale a:
\begin{enumerate}
    \item Un \textbf{tiempo de permanencia} (holding time) aleatorio antes de saltar, con distribución exponencial:
    \[ T_i \sim \text{Exp}(q_i) \implies E[T_i] = \frac{1}{-q_{ii}} \]
    
    \item Una \textbf{cadena de saltos embebida} que elige el destino $j$ con probabilidad:
    \[ \pi_{ij} = \frac{q_{ij}}{q_i} \quad \text{para } i \neq j \]
\end{enumerate}
\end{theorem}

\end{frame}



\section{Semigrupos y Exponencial Matricial}

\begin{frame}
\frametitle{Semigrupos de Transición}
\framesubtitle{Propiedades de la familia de matrices $P(t)$}

Sea $P(t)$ la matriz cuyas entradas son $p_{ij}(t) = P(X_{t+s}=j \mid X_s=i)$. En tiempo continuo, la familia de matrices $\{P(t)\}_{t \ge 0}$ forma un \textbf{semigrupo estándar de transición}.

\vspace{0.2cm}
Debe satisfacer tres propiedades fundamentales:



\end{frame}


\begin{frame}
\frametitle{Semigrupos de Transición}
\framesubtitle{El Generador como derivada en el origen}

¿Cómo se relaciona la matriz de probabilidades $P(t)$ con las tasas de salto $Q$? 

\vspace{0.3cm}
El generador infinitesimal $Q$ es, por definición, la \textbf{derivada del semigrupo} evaluada en $t = 0$:

\begin{theorem}[Generador del Semigrupo]
\[ Q = P'(0) = \lim_{h \to 0^+} \frac{P(h) - I}{h} \]
\end{theorem}

\vspace{0.2cm}
Esto implica que, para un intervalo de tiempo $h$ extremadamente pequeño, la matriz de transición se aproxima linealmente por las tasas de salto:
\[ P(h) = I + hQ + o(h) \]

\end{frame}

% fórmula de la Exponencial Matricial 
\begin{frame}
\frametitle{Solución Matricial}
\framesubtitle{La fórmula $P(t) = e^{tQ}$}

Al igual que la solución escalar de $y' = qy$ es $y(t) = e^{qt}$, la evolución del sistema en espacios de estados finitos se obtiene mediante la \textbf{exponencial de matrices}.

\begin{block}{Exponencial Matricial}
Se define a través de su serie de Taylor:
\[ P(t) = e^{tQ} = \sum_{n=0}^\infty \frac{t^n Q^n}{n!} = I + tQ + \frac{t^2 Q^2}{2!} + \dots \]
\end{block}
\end{frame}


\section{Ecuaciones Diferenciales de Kolmogorov}
\begin{frame}
\frametitle{Ecuaciones de Kolmogorov}
\framesubtitle{Derivación a partir de Chapman-Kolmogorov}

Para encontrar las probabilidades de transición, partimos de la propiedad fundamental del semigrupo (Chapman-Kolmogorov):
\[ P(t+h) = P(t)P(h) = P(h)P(t) \]

\end{frame}


\begin{frame}
\frametitle{Ecuaciones de Kolmogorov}
\framesubtitle{Ecuación Hacia Atrás }


\begin{block}{Kolmogorov}
\[ \mathbf{P'(t) = Q P(t)} \]
Condicionada sobre el \textbf{primer salto} que da el sistema al salir del estado inicial.
\end{block}

\vspace{0.2cm}
\footnotesize{\textit{Nota: Esta ecuación siempre es válida (incluso en espacios de estado infinitos) bajo condiciones mínimas de regularidad.}}
\end{frame}

% Ecuación Hacia Adelante
\begin{frame}
\frametitle{Ecuaciones de Kolmogorov}
\framesubtitle{Ecuación Hacia Adelante}

\begin{block}{Kolmogorov }
\[ \mathbf{P'(t) = P(t) Q} \]
Condicionada sobre el \textbf{último salto} que da el sistema antes del tiempo $t$.
\end{block}

\vspace{0.2cm}
\footnotesize{\textit{Nota: En espacios de estado infinitos, esta ecuación requiere condiciones más estrictas (como no explosión) para ser válida, ya que implica intercambiar límites y sumas infinitas.}}
\end{frame}



\section{Explosión y Regularidad}

\begin{frame}
\frametitle{Condición de Explosión}
\framesubtitle{¿Qué significa explotar en una CTMC?}


\begin{block}{Definición de Explosión}
Sea $T_n \sim \text{Exp}(-q_{X_n X_n})$ el tiempo que el sistema pasa en el $n$-ésimo estado visitado. El tiempo total para dar infinitos saltos es:
\[ \zeta = \sum_{n=0}^\infty T_n \]
Decimos que la cadena \textbf{explota} si $\zeta < \infty$ El sistema realiza infinitos saltos en un tiempo finito.
\end{block}

\vspace{0.2cm}
Si por el contrario, $P(\zeta = \infty) = 1$. En un sistema regular, la matriz $Q$

\end{frame}


\begin{frame}
\frametitle{Condición de Explosión}
\framesubtitle{El criterio exacto para Procesos de Nacimiento Puro}

Cuándo ocurre la explosión, analizamos un proceso de nacimiento puro (donde las transiciones solo van de $n \to n+1$) con tasas de nacimiento $\lambda_n$.

El tiempo total esperado para dar infinitos saltos es la suma de los tiempos esperados en cada estado:


\begin{theorem}[Condición de Explosión]
Un proceso de nacimiento puro con tasas $\lambda_n$ explota si y solo si la serie de los tiempos esperados converge:
\[ \sum_{n=0}^\infty \frac{1}{\lambda_n} < \infty \]
\end{theorem}

\end{frame}


\begin{frame}
\frametitle{Condición de Explosión}
\framesubtitle{Crecimiento lineal vs Cuadrático}

Comparemos dos sistemas de nacimiento puro usando el teorema anterior:

\vspace{0.3cm}
\textbf{1. Sistema Regular (Crecimiento Lineal): $\lambda_n = n$}
\[ \sum_{n=1}^\infty \frac{1}{n} =  \dots = \infty \]
Esta es la serie armónica, que diverge. \textbf{El sistema es regular}. Aunque los saltos son más rápidos, nunca alcanza el infinito en tiempo finito.

\vspace{0.3cm}
\textbf{2. Sistema Explosivo (Crecimiento Cuadrático): $\lambda_n = n^2$}
\[ \sum_{n=1}^\infty \frac{1}{n^2} = \dots = \frac{\pi^2}{6} \approx 1.645 \]
Esta serie converge.\textbf{El sistema explota} en un tiempo promedio de $\approx 1.645$.
\end{frame}

\end{document}